\documentclass[10pt,a4paper]{amsart}
\usepackage[margin=19mm]{geometry}
\usepackage{amsmath,amssymb,mathtools}
\usepackage[scr=rsfs,cal=euler]{mathalfa}
\usepackage{xfrac}
\usepackage[unicode,hidelinks,bookmarks=false]{hyperref}
\newcommand{\ie}{i.e.\ }
\newcommand{\cf}{cf.\ }
\newcommand{\resp}{respectively\ }
\newcommand{\Subsection}{\subsection}
\providecommand{\nobreakdash}{\nobreak\hbox{-}}
\setlength{\emergencystretch}{2em}
\multlinegap0pt
\usepackage{bm}
\usepackage{bbm}
\usepackage{dsfont}
\usepackage{xspace}
\usepackage{cancel}
\usepackage{mathtools}
\usepackage{amsthm}
\makeatletter
\@ifundefined{proposition}{\newtheorem{proposition}{Proposition}}{}
\makeatother
\title[Szeg\H{o}-Type Limit Theorem in the Drury-Arveson space]{Szeg\H{o}-Type Limit Theorem in the Drury-Arveson space}
\author{Arya Gayathri Memana}
\date{Source: arXiv:2508.00708v1 (CC BY 4.0)}
\begin{document}
\maketitle

\noindent\emph{Source and licence.} Szeg\H{o}-Type Limit Theorem in the Drury-Arveson space. Version arXiv:2508.00708v1 (CC BY 4.0), distributed under the Creative Commons Attribution 4.0 International licence, as recorded in the arXiv licence metadata for this exact version and shown on the abstract page https://arxiv.org/abs/2508.00708v1. Full rights notice: licenses/arxiv-2508.00708-CC-BY-4.0.txt.\par\smallskip

\noindent\textit{Editorial scope.} Counted unit: the Proposition whose source label is \texttt{None} in the article (source0.tex lines 445--460, source label \texttt{None}), delivered together with the article's own complete original proof of it (source0.tex lines 461--493, one \texttt{proof} environment of 33 lines). No mathematics is written, repaired or re-derived: statement and proof are the article's own text line for line. Required context retained uncounted: none. Every definition, display and label of those lines is retained unchanged. Omitted from this package and not redistributed: 1--444, 494--567. Nothing inside a retained excerpt is omitted or altered: every retained body is delivered exactly as the article's lines are written, so no omission receipt of any kind accompanies this package. Citations: the retained excerpts carry no citation command at all, so no bibliography entry of the article is transcribed and no reference is redistributed. Reference adaptation: none was needed and none was made; every reference inside the delivered unit resolves inside it, so no unresolved reference or citation is produced. Printed slips kept literal: no printed word, symbol, hypothesis or number of the counted statement or of its proof was altered. Layout and class: the article's own class is \texttt{amsart} with packages including babel, inputenc, amsmath, graphicx, enumerate, mathrsfs, mathtools, biblatex, tikz-cd, color, ulem, todonotes, enumitem, amsthm; the wrapper is the lane's pinned \texttt{amsart} editorial wrapper at 10pt on A4, which restates the theorem-like environments the retained text needs and adds the article's own macro definitions that the retained text needs (none), with extra wrapper packages (bm, bbm, dsfont, xspace, cancel, mathtools). The delivered numbering is the unit's own. Assets: no retained span contains a figure environment, a graphics-inclusion command, a PDF-page inclusion, a table or a TikZ picture; no image file is copied into this package, the package declares no asset dependency and no image file is redistributed with it. Licence: version arXiv:2508.00708v1 of this article is distributed under the Creative Commons Attribution 4.0 International licence, as recorded in the arXiv licence metadata for this exact version and shown on the abstract page https://arxiv.org/abs/2508.00708v1. A full rights notice is delivered at licenses/arxiv-2508.00708-CC-BY-4.0.txt. Characters: every retained excerpt is pure ASCII, so the delivered engine, XeLaTeX with its default Latin Modern text font, reports no missing character.
\begin{proposition}
   For a Toeplitz-like operator $T$ with a continuous symbol $\pi(T)=\phi$,
   the limit of the normalized trace exists, and 
   \begin{align*}
   \lim _{N \rightarrow \infty} \chi_N(T) = \int_{\partial \mathbb{B}_d} \phi(z,\bar{z})d\sigma
    \end{align*}
   %The tracial state $\chi$ on $C(\partial \mathbb{B}_d) \simeq \mathcal{T}/K(H_d^2)$ defined as
    %\begin{align*}
    %\chi(f) = \int_{\partial \mathbb{B}_d} f d\sigma
    %\end{align*}
    %where $d\sigma$ is the surface area measure on $\partial \mathbb{B}_d$ extends to the whole of $\mathcal{T}_d$ and this extension agrees with
    %\begin{align*}
     %   \chi(T) : = \lim_{n \rightarrow \infty} \chi_n(T).
    %\end{align*}
    %on $\mathcal{T}_d$.
\end{proposition}

\begin{proof}

Finite sums of the form 
\begin{align}\label{eq1}
 A = \sum_{i=1}^n p_i(S)^* q_i(S) + K
\end{align}
are dense in $\mathcal{T}_d$. For $A$, we have by the previous proposition,
\begin{align*}
\lim_{N \rightarrow \infty} \chi_N(A) = \int \sum_{i=1}^n \overline{p_i(z)} q_i(z) d\sigma
\end{align*}
Now for an arbitrary element $T$ in the $C^*$-algebra $\mathcal{T}_d$ and $\epsilon >0$ choose an element $A$ of the form as in $(\ref{eq1})$ such that $||T-A|| < \epsilon$. Choose $N'$ such that for all $M,N >N'$, $|\chi_M(A)- \chi_N(A)| < \epsilon$. Now we note that,
\begin{align*}
|\chi_N(T) - \chi_M(T)| & = |\chi_N(T) -\chi_N(A) + \chi_N(A) -\chi_M(A) +\chi_M(A) -\chi_M(T)| \\
& \leq |\chi_N (T)- \chi_N(A)| + |\chi_M(A) -\chi_N(A)| + |\chi_M(A) - \chi_M(T)|\\
& \leq ||\chi_N|| \ ||T-A|| + \epsilon + ||\chi_M|| \ ||T-A||\\
& \leq 2 ||T-A|| + \epsilon\\
& \leq 3 \epsilon
\end{align*}
for all $M,N > N'$. Thus, the sequence $\chi_N(T)$ converges for each $T \in \mathcal{T}_d$.\\ Let us call this limit $\chi(T)$. It follows that $\chi$ is also a state and $\chi$ is the unique $weak^*$ limit of $\chi_N$.\\

Since $\chi$ is zero on $K(H_d^2)$ it induces a state on the $C^*$-algebra $\mathcal{T}_d/K(H_d^2)$ which we continue to denote by $\chi$. We have proved so far that $\chi$ agrees with the linear functional $f \mapsto \int_{\partial \mathbb{B}_d}f d\sigma$ on $C(\partial \mathbb{B}_d)$ on the dense set of polynomials of the form ($\ref{eq1}$). Hence, by continuity, they agree everywhere and the result follows.


%Note that for an element $T_p +K $ of $\mathcal{T}_d/ K(H_d^2)$ with a polynomial symbol $p$, we have by the previous proposition that,
%\begin{align*}
%\lim_{n \rightarrow \infty} \chi_n(T_p+K) = \int_{\partial \mathbb{B}_d} p d \sigma
%\end{align*}
%Thus for $T_h \in \mathcal{T}_d+ K$ with continuous symbol $h$, we choose a sequence of polynomials $p_n$ such that $T_{p_n}+K$ converges to $T_h+K$ (this is because of the isomorchism in (2)). Thus we have,
%\begin{align*}
%\lim_{n \rightarrow \infty} \chi_n(T_h+K) = \int_{\partial \mathbb{B}_d} h d \sigma
%\end{align*}
%Since we know by Prop 1.1 $\chi(K) = 0$ for a compact operator
\end{proof}

\end{document}
